%!TEX program = uplatex
\RequirePackage{plautopatch}
\documentclass[uplatex,dvipdfmx]{jlreq}
\usepackage{amsmath,amssymb}

\pagestyle{empty}

\newcommand{\C}{\mathbb{C}}

\begin{document}

%======================================================================
% 表紙
%======================================================================
\begin{center}
  \vspace*{20pt}
  {\Huge\bfseries 剛性と対称性}\\[30pt]
  {\large いろいろな数学の集いの広場}\\[20pt]
  {\large 概要}
\end{center}

\vfill

\newpage

%======================================================================
% 講演１：金井雅彦「Mostow の剛性定理とその仲間達」
%======================================================================
\begin{center}
  {\large\bfseries Mostow の剛性定理とその仲間達}\\[4pt]
  金井 雅彦
\end{center}

\medskip

Lie 群の格子に関する剛性定理は種々の様相を有する。
数ある剛性定理の中でも多分最も著名であろう Mostow の剛性定理を例にとろう。
それによれば、
\begin{itemize}
  \item[(1)] 次元が 3 以上のコンパクト実双曲多様体は
             その基本群により位相・計量ともに完全に決定される。
\end{itemize}
これが Mostow の剛性定理の幾何学的定式化である。
また代数的な翻訳も瞬時にして可能である：
\begin{itemize}
  \item[(2)] 適当な仮定を満たす非コンパクト型半単純 Lie 群は
             その格子により完全に決定される。
\end{itemize}
いずれも、基本群ないし格子といった離散的対象により、一般にはより自由度の高い
はずの連続的な対象がコントロールされることを主張する。
さて、証明に目を転じてみよう。
Mostow 自身による証明（ちなみに現在では別証明が多数存在する）において
本質的な役割を果たしたのは、
\begin{itemize}
  \item[(3)] 擬等角写像の微分可能性
\end{itemize}
というある種解析的な事実であり、また
\begin{itemize}
  \item[(4)] 無限遠境界への基本群の作用に関するエルゴード理論的考察
\end{itemize}
であった。
ただし (4) に関しては、群作用により不変な共形構造の一意性が議論の中核をなすがゆえに、
より正しくいうならばエルゴード理論と幾何学のハイブリッドと言うべきかもしれない
（そう言えば、「Ergodic Geometry Seminar」というのがフランスにありましたが）。

さらに post-Mostow と言うことになれば、Margulis の超剛性定理などがすぐに思い浮かぶ。
特にこれにより高階の半単純 Lie 群の格子の算術性が保証される。
また Corlette、Gromov--Schoen、Jost--Yau、Mok--Siu--Yeung 等が
調和写像という大域変分法的な手法を用い Margulis の剛性定理の別証明を与えたのは
比較的記憶に新しい。
さらには群作用に対する剛性問題も忘れてはなるまい。
Mostow や Margulis の剛性定理においては格子から有限次元 Lie 群の中への準同型が
取り扱われたのに対し、群作用の剛性においては、群作用、
すなわち微分同相群という無限次元 Lie 群の中への準同型が対象となる。
この無限次元性から生じる困難をいかに克服するかがその醍醐味である。

このように剛性問題は種々の数学の交錯する場である。
この講演ではまず剛性問題を概観（第1話）した後に、
調和写像による超剛性定理の証明（第2話）、群作用の剛性（第3話）についてお話しする
予定である。

\newpage

%======================================================================
% 講演２：納谷信「モストウ以前 ── 揺藍期の剛性問題」
%======================================================================
\begin{center}
  {\large\bfseries モストウ以前 ── 揺藍期の剛性問題}
\end{center}
\begin{flushright}
  納谷 信（名古屋大学・多元数理）
\end{flushright}

\medskip

半単純リー群の格子の剛性の研究は、モストウの仕事によって大域的な剛性定理に到達し、
マルグリスの登場を待つこととなる。

この講演では、歴史を逆に遡り、「モストウ以前」を概観する。
とくに、ヴェイユの局所剛性定理に焦点を当て、この結果が
ある種のコホモロジー群の消滅に翻訳されることをみる。
その証明には、問題が調和積分論を用いてあざやかに解決されるさまを見ることができる。

\newpage

%======================================================================
% 講演３：井関裕靖「Mostow 剛性定理と最小エントロピー定理 I, II」
%======================================================================
\noindent
{\large\bfseries Mostow 剛性定理と最小エントロピー定理 I, II}
\hfill
{\large\bfseries 井関 裕靖（東北大・理）}

\medskip

Mostow 剛性定理の証明はいくつかありますが、
階数１の対称空間のココンパクト格子に対する証明は、
いずれも格子の同型から誘導される理想境界の間の同変な同相写像に注目し、
これが対称空間の間の同変な等長写像に拡張されることを示す、という手順を踏んでいます。
ここでは G.~Besson、G.~Courtois と S.~Gallot の三人による
重心写像を用いた証明のアイディアと概略を紹介します。
これは非常に簡明な証明で、彼らの論文にも
``easy enough to be taught eventually at the undergraduate level!''
と書かれています。
（とはいうものの、これを信じて4年生・大学院生向けの講義を始めた筆者は
後悔し始めています。もちろんこれはこの三人のせいでも証明のせいでもありません。）

この証明は最小エントロピー定理の証明の副産物として得られたものです。
この場合、エントロピーとは群作用の大きさを計る量で、
最小エントロピー定理は
「群作用がある意味で最小になるとき、その群作用はココンパクト格子の作用と一致する」
ということを主張しています。
これも広い意味での剛性定理だと言っていいでしょう。
ここでは関連したいくつかの話題の紹介をしたいと思っています。

\end{document}
